\chapter{Structs e Enumerações}
\label{cap:structs}

Até aqui, cada variável guardava um único valor de um único tipo. Esta seção apresenta ferramentas para agrupar dados relacionados em uma única unidade — essencial para modelar informações do mundo real, como a leitura de um sensor (valor, unidade, instante da medição) ou a configuração de um periférico.

\section{Registros e \code{structs}}

Uma \code{struct} agrupa múltiplas variáveis, de tipos possivelmente diferentes, sob um único nome — cada uma acessível por um \textbf{campo} (\emph{field}).

\begin{codigoc}{Definindo e usando uma struct}
#include <stdio.h>

typedef struct {
    char nome[20];
    float temperatura;
    int  umidade;
} Leitura;

int main(void) {
    Leitura sala = {"Sala de estar", 23.5f, 60};

    printf("%s: %.1f C, %d%% umidade\n",
           sala.nome, sala.temperatura, sala.umidade);

    sala.temperatura = 24.0f;   // acessando e alterando um campo
    return 0;
}
\end{codigoc}

O \code{typedef} antes de \code{struct} cria um \textbf{apelido de tipo}: em vez de escrever \code{struct Leitura sala;} toda vez, basta \code{Leitura sala;}. Campos são acessados com o operador ponto (\code{.}).

\subsection{Structs e ponteiros}

Quando se tem um \textbf{ponteiro} para uma struct (comum ao passá-la para funções, evitando copiar todos os seus campos), o acesso aos campos usa o operador seta (\code{->}) em vez do ponto:

\begin{codigoc}{Acessando campos via ponteiro}
#include <stdio.h>

typedef struct {
    float temperatura;
    int umidade;
} Leitura;

void exibir(Leitura *l) {
    printf("%.1f C, %d%%\n", l->temperatura, l->umidade);
    // equivalente a: (*l).temperatura, (*l).umidade
}

int main(void) {
    Leitura sala = {23.5f, 60};
    exibir(&sala);
    return 0;
}
\end{codigoc}

\section{Enumerações (\code{enum})}

Uma \code{enum} define um conjunto nomeado de constantes inteiras, tornando o código mais legível do que usar números ``mágicos'' soltos pelo texto.

\begin{codigoc}{Definindo uma enumeração}
#include <stdio.h>

typedef enum {
    ESTADO_AGUARDANDO,   // vale 0
    ESTADO_CONECTANDO,   // vale 1
    ESTADO_ATIVO,        // vale 2
    ESTADO_ERRO          // vale 3
} EstadoSistema;

int main(void) {
    EstadoSistema estado = ESTADO_CONECTANDO;

    if (estado == ESTADO_CONECTANDO) {
        printf("Conectando...\n");
    }

    return 0;
}
\end{codigoc}

Por padrão, o primeiro valor de uma \code{enum} vale 0, e cada valor seguinte é o anterior mais 1 — mas isso pode ser sobrescrito explicitamente (\code{ESTADO\_ERRO = 100}, por exemplo). Combinada com um \code{switch} (\cref{cap:decisao}), uma \code{enum} é a base clássica para implementar máquinas de estado em C.

\begin{esp32box}
\code{structs} e \code{enums} são ferramentas centrais para organizar o estado de um projeto de IoT. Uma \code{struct} agrupa naturalmente tudo o que pertence a uma mesma leitura de sensor — valor, unidade, instante da medição —, em vez de espalhar essa informação em várias variáveis soltas; uma \code{enum} representa, de forma legível, os diferentes \textbf{modos de operação} de um dispositivo (por exemplo, \code{MODO\_NORMAL}, \code{MODO\_ALERTA}, \code{MODO\_ECONOMIA}), testados com um \code{switch} a cada iteração do \code{loop()}. Usaremos exatamente essa combinação no projeto integrador do \cref{cap:projeto-final}.
\end{esp32box}
